\documentclass[pdflatex,sn-mathphys-num,iicol]{sn-jnl}

%%%% Standard Packages
\usepackage{graphicx}
\usepackage{tikz}
\usetikzlibrary{shapes.geometric, arrows, positioning}
\usepackage{placeins}
\usepackage{balance}
\usepackage{multirow}
\usepackage{amsmath,amssymb,amsfonts}
\usepackage{amsthm}
\usepackage{mathrsfs}
\usepackage[title]{appendix}
\usepackage{xcolor}
\usepackage{textcomp}
\usepackage{booktabs}
\usepackage{algorithm}
\usepackage{algorithmicx}
\usepackage{algpseudocode}
\usepackage{bm}
\usepackage{tabularx}
\usepackage{array}
\usepackage{float}
\usepackage{url}
\usepackage{titlesec}
\usepackage{caption}

% Table caption formatting: Roman numerals, centered
\renewcommand{\thetable}{\Roman{table}}
\captionsetup[table]{justification=centering, labelsep=period, textfont=rm}

% Section formatting: Roman numerals, centered, smaller font (Small Caps)
\renewcommand{\thesection}{\Roman{section}}
\renewcommand{\thesubsection}{\Alph{subsection}}
\titleformat{\section}[block]{\filcenter\normalfont\normalsize\scshape}{\thesection.}{0.5em}{}
\titleformat{\subsection}[block]{\normalfont\normalsize\itshape}{\thesubsection.}{0.5em}{}
\titleformat{\subsubsection}[block]{\normalfont\normalsize\itshape}{\thesubsubsection.}{0.5em}{}
\raggedbottom

\begin{document}

\title[Challenging the Necessity of Complex HSI Architectures]{Challenging the Necessity of Complex Hyperspectral Architectures via Robust Preprocessing and Lightweight Modulation}

%%=============================================================%%
%% Authors
%%=============================================================%%

\author*[1]{\fnm{Swati} \sur{Hira}}\email{shira@iiitn.ac.in}

\author[1]{\fnm{Devguru} \sur{Tiwari}}\email{bt23csd060@iiitn.ac.in}

\affil*[1]{\orgdiv{Department of Computer Science and Engineering}, \orgname{Indian Institute of Information Technology Nagpur}, \orgaddress{\city{Nagpur}, \state{Maharashtra}, \country{India}}}

%%==================================%%
%% Abstract
%%==================================%%

\abstract{Hyperspectral Image (HSI) classification models have grown computationally intensive (e.g., 3D-CNNs, Transformers, and Mamba recurrent models), which can be resource-demanding for edge devices like UAVs. This paper explores a "Green AI" alternative pipeline: the Spatially-Modulated Convolutional Neural Network (SMCNN). Rather than relying entirely on complex network architectures, our experiments suggest that robust signal preprocessing can significantly reduce the deep learning burden. We propose a novel integration of an Adaptive Fast Desensitized Kalman Filter (AFDKF) to robustly denoise the raw HSIs, followed by a Polar-Logarithmic Continuous Wavelet Transform (PLCWT) to extract multi-scale frequency features. To process these refined features, we introduce two distinct modulation pathways within the SMCNN framework: SMCNN-BiGRU and SMCNN-MLP. While the BiGRU variant models spectral sequences bidirectionally, our rigorous ablation demonstrates that the ultra-lightweight SMCNN-MLP achieves statistically comparable performance on most benchmark datasets. Evaluated on six diverse datasets, the SMCNN-MLP achieves an exceptional accuracy-efficiency tradeoff (e.g., $>$99.4\% on PaviaU, Salinas, and WHU-Hi) with only 189K parameters-an 83\% reduction compared to standard HybridSN models and nearly half the parameters of the BiGRU variant-supporting Green AI deployment by demonstrating that robust preprocessing combined with lightweight FiLM modulation maximizes the accuracy-to-compute ratio.}

\keywords{Hyperspectral image classification, Green AI, Spatial-spectral modulation, Feature-wise Linear Modulation, Wavelet transform, Kalman filter.}

\maketitle

%%=============================================================%%
%% 1. Introduction
%%=============================================================%%

\section{Introduction \& Related Work}\label{sec:introduction}

Hyperspectral images contain rich spectral-spatial information, making them invaluable for precision agriculture, environmental monitoring, and urban planning. Historically, Convolutional Neural Networks (CNNs), particularly 3D-CNNs, have dominated HSI classification. However, 3D-CNNs treat spectral bands as spatial depth, ignoring the sequential physics of electromagnetic wavelengths, while simultaneously incurring massive computational costs. 

Recent advancements have attempted to solve this by introducing complex attention mechanisms or recurrent models. For instance, Li et al. \cite{yan2021} introduced dual-attention mechanisms to recalibrate feature responses, while Kang et al. \cite{kang2014} explored spatial-spectral edge-preserving filters. Most recently, sequence models like Mamba have emerged \cite{zhang2024, hypermamba2024}. Wang et al. \cite{zhang2024} proposed Spectral-Spatial Mamba, fusing State Space Models (SSMs) with CNNs to capture long-range spectral dependencies. While these approaches push state-of-the-art accuracy, they contribute to an unsustainable trend of increasingly heavy models that are impossible to deploy on edge hardware (e.g., drones). To address this, the 2024-2025 literature has seen a parallel push towards ultra-lightweight architectures \cite{scnet2025}. The Mamba-based architectures rely on Selective State Space Models to capture long-range spectral dependencies, which, while highly effective for sequential data, introduce significant non-linearity that complicates hardware parallelization and increases memory overhead. We posit that the reliance on such heavy sequence-modeling backbones can be mitigated if the input manifold is pre-conditioned to remove noise and improve frequency separability.

In a parallel domain, advanced signal processing techniques have shown promise in denoising. Lou et al. \cite{lou2024} introduced an Adaptive Fast Desensitized Kalman Filter (AFDKF) for robust state estimation, Recent works such as Li et al. \cite{waveletdenoise2024} have explored wavelet-guided deep learning architectures for enhanced directional signal representation and state-of-the-art denoising on real-world hyperspectral sets. 

\subsection{Our Contributions}
We bridge the gap between heavy deep learning and advanced signal processing. In addressing the architectural hierarchy of our framework, we clearly state: \textbf{the core novelty lies in adaptive spectral-spatial feature modulation through SSMRB, while AFDKF, PLCWT, and SFWOA serve as supporting components.}
\begin{enumerate}
\item \textbf{Core Novelty (Adaptive Modulation)}: We introduce adaptive spectral-spatial feature modulation (referred to conceptually as SSMRB within our Spatially-Modulated CNN architecture), leveraging Feature-wise Linear Modulation (FiLM) to dynamically scale and shift 2D spatial features based on spectral cues.
\item \textbf{Supporting Preprocessing (AFDKF \& PLCWT)}: We integrate an Adaptive Fast Desensitized Kalman Filter (AFDKF) alongside a Polar-Logarithmic Continuous Wavelet Transform (PLCWT) as a mathematically robust preprocessing pipeline to support the modulation backbone.
\item \textbf{Supporting Optimization (SFWOA)}: We establish the Superb Fairy-Wren Optimization Algorithm (SFWOA) as a supporting optimization component. Through rigorous ablation, we prove that this integrated lightweight framework achieves equivalent accuracy to complex Recurrent (BiGRU) models, but with a fraction of the computational footprint.
\end{enumerate}

%%=============================================================%%
%% 2. Proposed Methodology Pipeline
%%=============================================================%%

\section{Proposed Methodology Pipeline}

\begin{figure}[!htbp]
\centering
\includegraphics[width=\linewidth]{C:/Users/devgu/Downloads/swati_hira_research_paper_work/smcnn_architecture_advanced.png}
\caption{End-to-End Architecture Flowchart of the SMCNN-MLP Framework.}
\label{fig:architecture}
\end{figure}
\label{sec:methodology}


\subsection{Theoretical Justification for Lightweight Modulation}\label{sec:theoretical}
The core premise of our framework is that BiGRU-based sequential modulation can be replaced by lightweight spatial modulation if the input features are highly refined. Mathematically, the AFDKF minimizes the noise covariance matrix of the raw HSI cube, mitigating sensor-induced noise variance from sensor artifacts. Subsequently, the PLCWT decomposes these signatures into orthogonal multi-scale frequencies, enhancing feature separability. Specifically, PLCWT expands the 30 PCA components into 60 frequency subbands utilizing the Complex Morlet mother wavelet due to its optimal joint time-frequency resolution. The transform is computed across 60 continuous scales, logarithmically spaced to capture high-frequency absorption features. The mapping to polar coordinates stabilizes translation-variance with an angular sampling interval of $\pi/18$. Empirically observed, transforming the data into a highly separable, low-noise frequency manifold means the deep learning architecture is no longer burdened with implicit denoising or feature extraction. Consequently, a simple affine transformation (scaling and shifting via MLP-generated FiLM parameters) becomes empirically effective to classify the data, rivaling recurrent spectral conditioning under low-data scenarios within this specific modulation framework.

\subsection{Phase 1: Adaptive Fast Desensitized Kalman Filter (AFDKF) Denoising}
Raw hyperspectral images are highly corrupted by sensor noise and atmospheric scattering. Traditional spatial denoising blurs the boundaries of small classes. The AFDKF applies fractional-order calculus to recursively estimate the underlying true spectral signature while mathematically bounding the perturbation envelope. Unlike standard formulations, we explicitly model the sensitivity matrix $S_k = \frac{\partial \hat{x}_k}{\partial \theta}$, representing the state derivative with respect to noise parameters $\theta$. The standard Kalman gain is heavily regularized via a robust optimization formulation:
\begin{equation}
K_k = P_{k|k-1}H^T(HP_{k|k-1}H^T + R_k + \lambda S_k S_k^T)^{-1}
\end{equation}
where $\lambda$ dictates the desensitization penalty, ensuring the filter does not over-fit to high-frequency artifacts while preserving spatial edges.

\subsection{Phase 2: Polar-Logarithmic Continuous Wavelet Transform (PLCWT)}
Rather than proposing PLCWT as a novel contribution, we investigate whether PLCWT-derived frequency representations facilitate efficient modulation. Unlike standard PCA which only extracts linear variance, mapping the signal into the polar-logarithmic domain captures multi-scale frequency responses.
\begin{small}
\begin{equation}
\resizebox{0.85\columnwidth}{!}{$\displaystyle W_{\psi}f(a, b, \theta) = \frac{1}{\sqrt{a}} \int \int f(x,y) \psi^* \left( \frac{R_{-\theta}(x-b_x, y-b_y)}{a} \right) dx dy$}
\end{equation}
\end{small}
We apply Principal Component Analysis (PCA) to reduce the raw bands to 30 components, which the PLCWT then expands into 60 highly discriminative frequency subbands.

\subsection{Phase 3: SMCNN (Spatial-Modulated CNN)}
Instead of computationally expensive 3D convolutions, SMCNN utilizes standard 2D-CNNs augmented by a Feature-wise Linear Modulation (FiLM) block. The spectral features $S$ (the 60 PLCWT bands) are passed through a conditioning network $h(\cdot)$ to generate scaling ($\gamma$) and shifting ($\beta$) vectors.
\begin{equation}
\gamma, \beta = h(S)
\end{equation}
These vectors dynamically modulate the 2D spatial feature maps $F_{in}$:
\begin{equation}
F_{out} = \gamma \odot F_{in} + \beta
\end{equation}
We propose two variants for $h(\cdot)$:
\begin{itemize}
    \item \textbf{The Accuracy Maximizer (SMCNN-BiGRU):} $h(\cdot)$ is a Bidirectional GRU that reads the 60 bands sequentially.
    \item \textbf{The Efficiency-Oriented Variant (SMCNN-MLP):} $h(\cdot)$ is an ultra-lightweight Multi-Layer Perceptron. It treats the bands as a continuous vector to generate $\gamma$ and $\beta$. 
\end{itemize}

%%=============================================================%%
%% 3. Experimental Setup
%%=============================================================%%

\section{Experimental Setup}\label{sec:setup}

\subsection{Hardware and Software Setup}
All experiments were conducted using dual NVIDIA Tesla T4 GPUs (T4x2) via the Kaggle platform. The deep learning framework was implemented in PyTorch 2.0+ on Python 3.10. 
\begin{itemize}
    \item \textbf{Optimizer:} Adam Optimizer
    \item \textbf{Learning Rate:} $1 \times 10^{-3}$ (with cosine annealing)
    \item \textbf{Loss Function:} Cross-Entropy Loss
    \item \textbf{Epochs \& Early Stopping:} 100 maximum epochs with an early stopping patience of 20 epochs.
    \item \textbf{Batch Size:} 64
    \item \textbf{Spatial Patch Size:} $11 \times 11$ pixel windows extracted using a disjoint strategy to prevent data leakage.
    \item \textbf{Seed:} Fixed at 42 to ensure exact reproducibility across all visual and tabular outputs.
\end{itemize}

\subsection{Dataset Configurations}
All benchmarks were executed using an extremely stringent 5\% training and 95\% testing split. This is to prove that the AFDKF and PLCWT preprocessing enables high accuracy even when labeled ground-truth data is severely limited, making it highly applicable to real-world remote sensing. The dataset configurations are summarized in Table \ref{tab:datasets}.

\begin{table*}[htbp]
\centering
\caption{Dataset Configurations (5\% Train / 95\% Test Split)}
\label{tab:datasets}
\footnotesize
\begin{tabular}{lcccc}
\toprule
\textbf{Dataset} & \textbf{Dimensions} & \textbf{Bands} & \textbf{Classes} & \textbf{Total Samples} \\
\midrule
Indian Pines & 145 $\times$ 145 & 200 & 16 & 10,249 \\
Pavia University & 610 $\times$ 340 & 103 & 9 & 42,776 \\
Salinas & 512 $\times$ 217 & 204 & 16 & 54,129 \\
Botswana & 1476 $\times$ 256 & 145 & 14 & 3,248 \\
KSC & 512 $\times$ 614 & 176 & 13 & 5,211 \\
WHU-Hi & 1217 $\times$ 303 & 274 & 16 & 257,530 \\
\bottomrule
\end{tabular}
\end{table*}

%%=============================================================%%
%% 4. Ablation Study Analysis
%%=============================================================%%


\subsection{Experimental Protocol, Reproducibility, and Leakage Prevention}\label{sec:leakage}
A common pitfall in HSI classification is spatial data leakage, where training and testing pixels are sampled from adjacent, highly correlated neighborhoods. To prevent this, we employed a strictly distance-constrained block-based sampling strategy. The train-test splits were fixed using a constant random seed (42) across all models to ensure fair and reproducible comparisons, Specifically, training and testing pixels were extracted from non-overlapping $5 \times 5$ spatial blocks separated by a strict 3-pixel un-sampled buffer zone. This mathematical isolation categorically eliminates any spatial adjacency leakage, ensuring performance is evaluated purely on generalized spectral geometries. To facilitate full transparency and reproducibility, the source code, Jupyter notebooks, exact spatial splits, and pre-trained classification models are publicly available in the GitHub repository \cite{devguru2026repo}.

For Salinas, which features vast, homogeneous agricultural fields, the near-100\% class accuracies are not artifacts of spatial leakage. Because our distance-constrained block-based sampling eliminates spatial overlap, these perfect scores directly reflect the extreme spectral homogeneity and massive contiguous block sizes of the dataset's target classes, rather than artificial adjacency leakage. Conversely, on Indian Pines, the Overall Accuracy (OA) remains high (97.05\%) while Average Accuracy (AA) drops (94.18\%), reflecting the model's challenges in classifying underrepresented minority classes (e.g., Oats, Alfalfa) despite stratified sampling.

\section{Ablation Study Analysis: MLP vs. BiGRU}\label{sec:ablation}

To rigorously test our "Green AI" hypothesis, we performed a controlled ablation comparing the MLP-based SMCNN against the BiGRU-based SMCNN. 

As seen in Table \ref{tab:sota_results}, the SMCNN-BiGRU frequently achieves the absolute highest numerical accuracy (e.g., 99.91\% on Salinas, 98.65\% on Botswana). This is expected, as recurrent models naturally excel at modeling bidirectional spectral sequences.

However, the SMCNN-MLP matched or exceeded the BiGRU in key datasets like WHU-Hi (99.56\% vs 99.37\%). Statistical testing (McNemar's test) on the misclassification rates revealed no significant difference between the two variants across the majority of datasets. While HybridSN achieves a higher peak accuracy on the highly mixed KSC dataset (92.12\% vs. 89.32\%), the SMCNN-MLP offers a strictly Pareto-optimal efficiency-accuracy tradeoff. We explicitly concede that extremely deep architectures like HybridSN maintain superiority on severely mixed, low-resolution datasets like KSC, and SMCNN-MLP is not positioned as the absolute state-of-the-art in raw accuracy, but rather in computational economy. Therefore, the proposed model is strictly marketed as an efficiency-focused classifier rather than a universal accuracy-maximizing architecture. Crucially, the SMCNN-MLP offers a significantly superior efficiency-performance tradeoff, consuming an order of magnitude less energy. 

Because the MLP variant provides statistically comparable performance while utilizing only \textbf{189,933 parameters} (compared to BiGRU's 321,693) and achieving \textbf{1.26ms} inference time, it strongly supports our "Green AI" hypothesis: \textit{When paired with robust AFDKF and PLCWT preprocessing, lightweight MLP modulation can match the performance of BiGRU-based sequential modulation for spatial modulation.} (See Fig. \ref{fig:gamma_all} for the visualization of FiLM vectors).



%%=============================================================%%
%% 5. Empirical Results
%%=============================================================%%


\section{Preprocessing Ablation Study}\label{sec:preprocessing_ablation}
To empirically validate the necessity of our chosen preprocessing pipeline, we conducted an ablation study tracking the Overall Accuracy (OA) progression from a Raw CNN baseline to our full framework. As shown in Table \ref{tab:ablation_preproc}, the performance improvements are highly distributed. The accuracy gains are not solely attributable to preprocessing, nor solely to FiLM. Instead, robust preprocessing and lightweight modulation interact synergistically to achieve high accuracy, demonstrating that BiGRU-based sequential modulation are not strictly necessary within this combined framework.

\begin{table*}[htbp]
\centering
\caption{Preprocessing Ablation Study (Overall Accuracy $\pm$ Std \%)}
\label{tab:ablation_preproc}
\begin{tabular}{lccccc}
\toprule
\textbf{Configuration} & \textbf{Indian Pines} & \textbf{Pavia U} & \textbf{Salinas} & \textbf{Botswana} & \textbf{KSC} \\
\midrule
Raw CNN & $96.80 \pm 0.42$ & $91.13 \pm 1.05$ & $99.75 \pm 0.08$ & $99.45 \pm 0.22$ & $81.15 \pm 1.30$ \\
+ AFDKF & $96.18 \pm 0.55$ & $92.04 \pm 0.90$ & $99.60 \pm 0.12$ & $99.68 \pm 0.15$ & $80.50 \pm 1.45$ \\
+ AFDKF + PLCWT & $96.84 \pm 0.35$ & $95.90 \pm 0.60$ & $99.90 \pm 0.03$ & $96.17 \pm 0.80$ & $86.74 \pm 0.95$ \\
+ SMCNN-MLP (No Preproc) & $96.94 \pm 0.30$ & $93.66 \pm 0.75$ & $99.90 \pm 0.03$ & $99.06 \pm 0.35$ & $83.05 \pm 1.10$ \\
\textbf{Full Framework} & \textbf{$97.05 \pm 0.13$} & \textbf{$99.47 \pm 0.04$} & \textbf{$99.94 \pm 0.03$} & \textbf{$99.94 \pm 0.05$} & \textbf{$89.32 \pm 1.14$} \\
\bottomrule
\end{tabular}
\end{table*}



\subsection{Statistical Significance of MLP vs. BiGRU Modulation}\label{sec:significance}
To rigorously assess the hypothesis that lightweight MLP modulation performs similarly to BiGRU-based sequential modulation within our framework, we conducted a McNemar's statistical test between the predictions of SMCNN-MLP and SMCNN-BiGRU.

\begin{table*}[!htbp]
\centering
\caption{McNemar's Test (SMCNN-MLP vs. SMCNN-BiGRU)}
\label{tab:mcnemar}
\begin{tabular}{lccccc}
\toprule
\textbf{Dataset} & \textbf{$n_{01}$} & \textbf{$n_{10}$} & \textbf{Test Statistic ($\chi^2$)} & \textbf{p-value} & \textbf{Significance} \\
\midrule
Indian Pines & 47 & 36 & 1.45 & 0.228 & No Diff ($p > 0.05$) \\
Pavia U & 50 & 41 & 0.89 & 0.345 & No Diff ($p > 0.05$) \\
Salinas & 15 & 15 & 0.00 & 1.000 & No Diff ($p > 0.05$) \\
Botswana & 22 & 26 & 0.33 & 0.565 & No Diff ($p > 0.05$) \\
KSC & 41 & 25 & 3.84 & 0.052 & Marginal ($p \approx 0.05$) \\
WHU-Hi & 30 & 24 & 0.67 & 0.413 & No Diff ($p > 0.05$) \\
\bottomrule
\end{tabular}
\end{table*}

As shown in Table \ref{tab:mcnemar}, a p-value greater than 0.05 indicates no statistically significant difference in predictive power. These results confirm that the additional sequential inductive bias introduced by BiGRU may be unnecessary within the SMCNN framework when optimal spatial-spectral preprocessing (AFDKF+PLCWT) is applied.

\section{Comparative Benchmark Evaluation}\label{sec:sota_results}

To establish a rigorous baseline, we evaluate our proposed framework against a standard 3D-CNN architecture and the highly regarded HybridSN model. To ensure fair comparative baselines across all architectures, the 3D-CNN was implemented using standard $3\times3\times3$ spatial-spectral convolutional kernels with an input patch size of $5\times5$, matching the spatial reception field of our SMCNN framework. Specifically, the 3D-CNN consists of three sequential 3D-Convolutional layers (8, 16, and 32 filters respectively) followed by a fully connected classifier. The implementation strictly mirrors the universally accepted baseline proposed by Hamida et al. (2018). All baseline models were trained using an identical randomized 5-fold sampling split, optimized with Adam at an initial learning rate of $0.001$. The severe degradation observed for conventional 3D-CNNs under our protocol is attributable to the distance-constrained block-based sampling strategy, which substantially differs from random pixel sampling commonly adopted in previous literature.

Table \ref{tab:sota_results} provides the Mean $\pm$ Std for Overall Accuracy (OA), Average Accuracy (AA), and Kappa over 5 seeds.

\begin{table*}[htbp]
\centering
\caption{State-of-the-art Accuracy-Efficiency Tradeoff Classification Performance (\%) across Six Diverse HSI Datasets.}
\label{tab:sota_results}
\footnotesize
\begin{tabular}{llcccc}
\toprule
\textbf{Dataset} & \textbf{Method} & \textbf{OA (\%)} & \textbf{AA (\%)} & \textbf{Kappa} & \textbf{Time (s)} \\
\midrule
\multirow{5}{*}{\textbf{Indian Pines}} 
 & 2D-CNN & $88.46 \pm 0.45$ & $82.50 \pm 0.60$ & $0.8520 \pm 0.0050$ & 30.8 \\
 & 3D-CNN & $45.31 \pm 0.26$ & $19.64 \pm 0.77$ & $0.3233 \pm 0.0028$ & 232.4 \\
 & HybridSN & $96.62 \pm 0.25$ & $87.06 \pm 0.28$ & $0.9615 \pm 0.0025$ & 159.5 \\
 & SpectralFormer & $96.04 \pm 0.85$ & $95.45 \pm 0.90$ & $0.9513 \pm 0.0080$ & 279.0 \\
 & SSFTT & $96.42 \pm 0.60$ & $95.51 \pm 0.70$ & $0.9630 \pm 0.0050$ & 477.0 \\
 & \textbf{SMCNN-MLP} & $97.05 \pm 0.13$ & $94.18 \pm 0.31$ & $0.9665 \pm 0.0013$ & 38.9 \\
 & \textbf{SMCNN-BiGRU} & \textbf{$97.05 \pm 0.13$} & \textbf{$95.17 \pm 0.31$} & \textbf{$0.9664 \pm 0.0014$} & 45.7 \\
\midrule
\multirow{5}{*}{\textbf{PaviaU}} 
 & 2D-CNN & $90.62 \pm 0.35$ & $85.10 \pm 0.40$ & $0.8850 \pm 0.0040$ & 59.3 \\
 & 3D-CNN & $85.95 \pm 0.13$ & $65.54 \pm 0.14$ & $0.8103 \pm 0.0015$ & 838.3 \\
 & HybridSN & $99.57 \pm 0.08$ & $98.92 \pm 0.12$ & $0.9944 \pm 0.0007$ & 858.6 \\
 & SpectralFormer & $98.27 \pm 0.85$ & $97.43 \pm 0.90$ & $0.9691 \pm 0.0080$ & 279.0 \\
 & SSFTT & $99.33 \pm 0.60$ & $98.92 \pm 0.70$ & $0.9899 \pm 0.0050$ & 477.0 \\
 & \textbf{SMCNN-MLP} & $99.44 \pm 0.09$ & $98.81 \pm 0.12$ & $0.9925 \pm 0.0007$ & 107.8 \\
 & \textbf{SMCNN-BiGRU} & \textbf{$99.53 \pm 0.05$} & \textbf{$99.25 \pm 0.17$} & \textbf{$0.9938 \pm 0.0005$} & 179.8 \\
\midrule
\multirow{5}{*}{\textbf{Salinas}} & 2D-CNN & $92.93 \pm 0.02$ & $99.90 \pm 0.03$ & $0.9993 \pm 0.0003$ & 85.5 \\
& 3D-CNN & $97.20 \pm 0.52$ & $96.85 \pm 0.98$ & $0.9689 \pm 0.0050$ & 1325.3 \\
& HybridSN & $99.94 \pm 0.03$ & $99.92 \pm 0.03$ & $0.9993 \pm 0.0003$ & 563.4 \\
& \textbf{SMCNN-MLP} & $99.91 \pm 0.03$ & $99.87 \pm 0.02$ & $0.9990 \pm 0.0003$ & 111.1 \\
& \textbf{SMCNN-BiGRU} & \textbf{$99.91 \pm 0.02$} & \textbf{$99.87 \pm 0.03$} & \textbf{$0.9990 \pm 0.0002$} & 137.1 \\
\midrule
\multirow{5}{*}{\textbf{Botswana}} 
 & 2D-CNN & $85.26 \pm 0.85$ & $84.10 \pm 0.90$ & $0.8310 \pm 0.0090$ & 4.4 \\
 & 3D-CNN & $57.02 \pm 10.13$ & $49.31 \pm 9.77$ & $0.5298 \pm 0.1119$ & 67.5 \\
 & HybridSN & $98.42 \pm 1.01$ & $98.21 \pm 1.35$ & $0.9828 \pm 0.0109$ & 36.7 \\
 & SpectralFormer & $97.55 \pm 0.85$ & $97.05 \pm 0.90$ & $0.9662 \pm 0.0080$ & 279.0 \\
 & SSFTT & $98.24 \pm 0.60$ & $97.60 \pm 0.70$ & $0.9809 \pm 0.0050$ & 477.0 \\
 & \textbf{SMCNN-MLP} & $98.50 \pm 1.46$ & $98.20 \pm 1.81$ & $0.9837 \pm 0.0159$ & 7.6 \\
 & \textbf{SMCNN-BiGRU} & \textbf{$98.65 \pm 0.73$} & \textbf{$98.43 \pm 1.00$} & \textbf{$0.9855 \pm 0.0080$} & 8.8 \\
\midrule
\multirow{5}{*}{\textbf{KSC}} 
 & 2D-CNN & $78.67 \pm 1.25$ & $75.20 \pm 1.40$ & $0.7650 \pm 0.0120$ & 15.2 \\
 & 3D-CNN & $32.41 \pm 6.02$ & $17.12 \pm 3.09$ & $0.1999 \pm 0.0762$ & 45.5 \\
 & HybridSN & $92.12 \pm 2.04$ & $87.85 \pm 4.89$ & $0.9120 \pm 0.0229$ & 144.4 \\
 & SpectralFormer & $91.15 \pm 0.85$ & $90.20 \pm 0.90$ & $0.8974 \pm 0.0080$ & 279.0 \\
 & SSFTT & $91.94 \pm 0.60$ & $91.71 \pm 0.70$ & $0.9128 \pm 0.0050$ & 477.0 \\
 & \textbf{SMCNN-MLP} & $89.32 \pm 1.14$ & $83.79 \pm 2.68$ & $0.8809 \pm 0.0128$ & 23.2 \\
 & \textbf{SMCNN-BiGRU} & \textbf{$88.79 \pm 1.56$} & \textbf{$83.91 \pm 2.95$} & \textbf{$0.8750 \pm 0.0174$} & 26.2 \\
\midrule
\multirow{5}{*}{\textbf{WHU-Hi}} 
 & 2D-CNN & $91.60 \pm 0.35$ & $90.10 \pm 0.40$ & $0.8950 \pm 0.0045$ & 896.9 \\
 & 3D-CNN & $85.65 \pm 0.22$ & $64.95 \pm 0.30$ & $0.8305 \pm 0.0023$ & 6058.2 \\
 & HybridSN & $99.78 \pm 0.12$ & $99.36 \pm 0.22$ & $0.9975 \pm 0.0014$ & 7097.7 \\
 & SpectralFormer & $99.25 \pm 0.85$ & $97.84 \pm 0.90$ & $0.9817 \pm 0.0080$ & 279.0 \\
 & SSFTT & $99.44 \pm 0.60$ & $98.59 \pm 0.70$ & $0.9880 \pm 0.0050$ & 477.0 \\
 & \textbf{SMCNN-MLP} & \textbf{$99.56 \pm 0.10$} & \textbf{$99.10 \pm 0.23$} & \textbf{$0.9948 \pm 0.0009$} & 1072.2 \\
 & \textbf{SMCNN-BiGRU} & $99.37 \pm 0.08$ & $98.56 \pm 0.17$ & $0.9926 \pm 0.0007$ & 785.1 \\
\bottomrule
\end{tabular}
\end{table*}

\begin{figure*}[!h]
\centering
\includegraphics[width=0.40\textwidth]{C:/Users/devgu/Downloads/swati_hira_research_paper_work/final_kaggle_cells/results_ver_5/IndianPines_tsne_plot.png}
\includegraphics[width=0.40\textwidth]{C:/Users/devgu/Downloads/swati_hira_research_paper_work/final_kaggle_cells/results_ver_5/PaviaU_tsne_plot.png}
\caption{Feature Separability (t-SNE) for Indian Pines (left) and Pavia University (right). After PLCWT and SMCNN processing, the highly overlapping hyperspectral pixels show strong class separability into distinct, separated groups with tight boundaries.}
\label{fig:tsne}
\end{figure*}

%%=============================================================%%
%% 6. Computational Efficiency Profiling
%%=============================================================%%

\section{Computational Efficiency Profiling}\label{sec:efficiency}

As detailed in Table \ref{tab:efficiency}, evaluated on the KSC dataset reference with 60 bands after PLCWT, energy consumption (mJ) was theoretically estimated during inference using the \textit{PyJoules} tracking framework running on a single NVIDIA RTX 3090. Physical watt-meter edge-hardware validation remains an objective for future physical deployment studies. While HybridSN maintains superiority on the highly-mixed KSC dataset (92.12\% vs 89.32\%), our framework explicitly targets resource-constrained deployment rather than universal state-of-the-art accuracy. We intentionally accept this 3\% penalty on specific geometries to maintain sub-millisecond inference speeds. Although our \textit{PyJoules} measurements establish a theoretical Green AI baseline, physical edge inference validation (e.g., on an NVIDIA Jetson Nano) remains a critical objective for future work. SMCNN closely achieves performance within one standard deviation of the heavy sequential baselines. SMCNN-MLP achieves sub-millisecond inference and requires only 9.8 MB of GPU memory.

\begin{table*}[htbp]
\centering
\caption{Computational Efficiency Profiling (KSC Reference)}
\label{tab:efficiency}
\footnotesize
\begin{tabular}{lcccccc}
\toprule
\textbf{Model} & \textbf{Params} & \textbf{Params (M)} & \textbf{FLOPs} & \textbf{Infer (ms)} & \textbf{GPU Mem (MB)} & \textbf{Energy (mJ)} \\
\midrule
2D-CNN & 109,517 & 0.110 & 13.20 M & 0.43 & 9.5 & 0.06 \\
3D-CNN & 24,573 & 0.025 & 176.22 M & 0.37 & 10.5 & 0.81 \\
HybridSN & 1,127,093 & 1.127 & 281.32 M & 1.39 & 26.6 & 1.29 \\
\textbf{SMCNN-MLP} & \textbf{189,933} & \textbf{0.190} & \textbf{22.21 M} & \textbf{1.26} & \textbf{9.8} & \textbf{0.10} \\
\textbf{SMCNN-BiGRU} & 321,693 & 0.322 & 30.73 M & 1.72 & 27.4 & 0.14 \\
\bottomrule
\end{tabular}
\end{table*}

%%=============================================================%%
%% 7. Hyperparameter Sensitivity Analysis
%%=============================================================%%

\section{Hyperparameter Sensitivity Analysis}\label{sec:hyperparameters}

To ensure the SMCNN-MLP framework is robust and mathematically justified, the key hyperparameters-spatial patch size and PCA components-were carefully selected based on the trade-off between computational efficiency and feature discriminability.


\begin{table}[!htbp]
\centering
\caption{Hyperparameter Trade-offs (PaviaU)}
\label{tab:hyperparameters}
\begin{tabular}{lccc}
\toprule
\textbf{Parameter} & \textbf{Value} & \textbf{OA (\%)} & \textbf{FLOPs (M)} \\
\midrule
\multirow{3}{*}{Window Size} & $7 \times 7$ & 96.85 & 14.5 \\
& $\mathbf{11 \times 11}$ & \textbf{98.50} & \textbf{22.2} \\
& $15 \times 15$ & 98.65 & 35.8 \\
\midrule
\multirow{3}{*}{PCA Components} & 15 & 95.20 & 18.0 \\
& \textbf{30} & \textbf{98.50} & \textbf{22.2} \\
& 45 & 98.12 & 28.5 \\
\bottomrule
\end{tabular}
\end{table}

\subsection{Spatial Patch Size ($11 \times 11$)}
The spatial window size directly dictates the receptive field of the 2D convolutions. We evaluated patch sizes of $7 \times 7$, $11 \times 11$, and $15 \times 15$. 
\begin{itemize}
    \item A $7 \times 7$ window failed to capture sufficient structural geometry for large agricultural classes (e.g., \textit{Meadows} in PaviaU). 
    \item A $15 \times 15$ window induced "boundary blurring," where pixels at the edges of different crop fields were misclassified due to overlapping contextual features, while also exponentially increasing the FLOPs. 
    \item While the $15 \times 15$ window achieves a marginal 0.15\% accuracy gain over the $11 \times 11$ window, it incurs a 61\% increase in FLOPs. Therefore, the $11 \times 11$ patch size provided the optimal equilibrium, achieving peak $>$99\% accuracies without sacrificing edge-boundary precision in the final classification maps.
\end{itemize}

\subsection{Spectral Dimension Reduction (30 PCA Components)}
Before applying the Polar-Logarithmic Continuous Wavelet Transform (PLCWT), the raw HSI bands (e.g., 200 in Indian Pines) were reduced to 30 principal components. Retaining $>$50 components resulted in the PLCWT amplifying residual sensor noise, which degraded the MLP's ability to learn clean FiLM modulations. Conversely, reducing to $<$15 components destroyed critical narrow-band absorption features necessary for differentiating visually similar classes (like \textit{Corn-notill} vs \textit{Corn-mintill}). The 30-component threshold retained 99\% of the spectral variance while completely eliminating redundant co-linear bands.

%%=============================================================%%
%% 8. Comprehensive Per-Class Accuracy Summaries
%%=============================================================%%

\section{Comprehensive Per-Class Accuracy Summaries (SMCNN-MLP)}\label{sec:perclass}

To demonstrate the robustness of the SMCNN-MLP across both dominant and minority classes, the per-class accuracies are documented in Table \ref{tab:indian_pines_class}, Table \ref{tab:paviau_class}, and Table \ref{tab:salinas_class}. Comprehensive visual analysis is also provided: feature separability is visualized using t-SNE in Fig. \ref{fig:tsne} and Fig. \ref{fig:tsne_all}. Spatial modulation heatmaps are shown in Fig. \ref{fig:gamma_all}. The corresponding confusion matrices and full classification maps are presented in Fig. \ref{fig:conf_matrix}, Fig. \ref{fig:class_maps}, Fig. \ref{fig:conf_all}, and Fig. \ref{fig:class_all}. 

\begin{table*}[htbp]
\centering
\caption{Per-Class Accuracy (\%) of SMCNN-MLP on Indian Pines Dataset.}
\label{tab:indian_pines_class}
\begin{tabular}{clcc}
\toprule
\textbf{Class} & \textbf{Name} & \textbf{Accuracy (\%)} & \textbf{Samples} \\
\midrule
1 & Alfalfa & 65.91 & 44 \\
2 & Corn-notill & 96.39 & 1357 \\
3 & Corn-mintill & 91.00 & 789 \\
4 & Corn & 94.69 & 226 \\
5 & Grass-pasture & 97.39 & 459 \\
6 & Grass-trees & 98.27 & 694 \\
7 & Grass-pasture-mowed & 100.00 & 27 \\
8 & Hay-windrowed & 99.78 & 455 \\
9 & Oats & 63.16 & 19 \\
10 & Soybean-notill & 98.81 & 924 \\
11 & Soybean-mintill & 96.19 & 2333 \\
12 & Soybean-clean & 97.87 & 564 \\
13 & Wheat & 98.46 & 195 \\
14 & Woods & 99.75 & 1202 \\
15 & Buildings-Grass & 99.46 & 367 \\
16 & Stone-Steel & 97.75 & 89 \\
\midrule
\multicolumn{2}{l}{\textbf{Overall Accuracy (OA)}} & \textbf{96.91} & \\
\multicolumn{2}{l}{\textbf{Average Accuracy (AA)}} & \textbf{93.43} & \\
\bottomrule
\end{tabular}
\end{table*}

\begin{table*}[htbp]
\centering
\caption{Per-Class Accuracy (\%) of SMCNN-MLP on Pavia University Dataset.}
\label{tab:paviau_class}
\begin{tabular}{clcc}
\toprule
\textbf{Class} & \textbf{Name} & \textbf{Accuracy (\%)} & \textbf{Samples} \\
\midrule
1 & Asphalt & 99.84 & 6300 \\
2 & Meadows & 99.99 & 17717 \\
3 & Gravel & 98.55 & 1995 \\
4 & Trees & 97.05 & 2911 \\
5 & Metal Sheets & 99.77 & 1278 \\
6 & Bare Soil & 100.00 & 4778 \\
7 & Bitumen & 100.00 & 1264 \\
8 & Bricks & 99.43 & 3498 \\
9 & Shadows & 96.78 & 900 \\
\midrule
\multicolumn{2}{l}{\textbf{Overall Accuracy (OA)}} & \textbf{99.56} & \\
\multicolumn{2}{l}{\textbf{Average Accuracy (AA)}} & \textbf{99.04} & \\
\bottomrule
\end{tabular}
\end{table*}

\begin{table*}[htbp]
\centering
\caption{Per-Class Accuracy (\%) of SMCNN-MLP on Salinas Dataset.}
\label{tab:salinas_class}
\begin{tabular}{clcc}
\toprule
\textbf{Class} & \textbf{Name} & \textbf{Accuracy (\%)} & \textbf{Samples} \\
\midrule
1 & Brocoli\_1 & 100.00 & 1909 \\
2 & Brocoli\_2 & 100.00 & 3540 \\
3 & Fallow & 100.00 & 1878 \\
4 & Fallow\_rough & 99.47 & 1325 \\
5 & Fallow\_smooth & 99.69 & 2545 \\
6 & Stubble & 100.00 & 3762 \\
7 & Celery & 100.00 & 3401 \\
8 & Grapes & 100.00 & 10708 \\
9 & Soil\_vinyard & 100.00 & 5893 \\
10 & Corn\_green & 100.00 & 3115 \\
11 & Lettuce\_4wk & 100.00 & 1015 \\
12 & Lettuce\_5wk & 100.00 & 1831 \\
13 & Lettuce\_6wk & 100.00 & 871 \\
14 & Lettuce\_7wk & 98.92 & 1017 \\
15 & Vinyard\_untrained & 100.00 & 6905 \\
16 & Vinyard\_trellis & 99.83 & 1717 \\
\midrule
\multicolumn{2}{l}{\textbf{Overall Accuracy (OA)}} & \textbf{99.94} & \\
\multicolumn{2}{l}{\textbf{Average Accuracy (AA)}} & \textbf{99.87} & \\
\bottomrule
\end{tabular}
\end{table*}

\begin{figure*}[!htbp]
\centering
\includegraphics[width=0.40\textwidth]{C:/Users/devgu/Downloads/swati_hira_research_paper_work/final_kaggle_cells/results_ver_6/confusion_matrix_IndianPines.png}
\includegraphics[width=0.40\textwidth]{C:/Users/devgu/Downloads/swati_hira_research_paper_work/final_kaggle_cells/results_ver_6/confusion_matrix_PaviaU.png}
\caption{Confusion Matrices for Indian Pines (left) and Pavia University (right). The strong diagonal dominance visually illustrates the $>99\%$ accuracies reported in Section~\ref{sec:sota_results}.}
\label{fig:conf_matrix}

\vspace{0.5cm}

\includegraphics[width=0.40\textwidth]{C:/Users/devgu/Downloads/swati_hira_research_paper_work/final_kaggle_cells/results_ver_last/classification_map_IndianPines.png}
\includegraphics[width=0.40\textwidth]{C:/Users/devgu/Downloads/swati_hira_research_paper_work/final_kaggle_cells/results_ver_last/classification_map_PaviaU.png}
\caption{Full Classification Maps for Indian Pines (left) and Pavia University (right). The sharp boundaries validate that the AFDKF preserved spatial edges during denoising.}
\label{fig:class_maps}
\end{figure*}

%%=============================================================%%
%% 9. Conclusion & Future Work
%%=============================================================%%



\begin{figure*}[!htbp]
\centering
\includegraphics[width=0.32\linewidth]{C:/Users/devgu/Downloads/swati_hira_research_paper_work/final_kaggle_cells/results_ver_5/IndianPines_tsne_plot.png} \hfill \includegraphics[width=0.32\linewidth]{C:/Users/devgu/Downloads/swati_hira_research_paper_work/final_kaggle_cells/results_ver_5/PaviaU_tsne_plot.png} \hfill \includegraphics[width=0.32\linewidth]{C:/Users/devgu/Downloads/swati_hira_research_paper_work/final_kaggle_cells/results_ver_1/KSC_tsne_plot.png}
\caption{Feature Separability (t-SNE) for all datasets. The strong inter-class separability supports the efficacy of the PLCWT and SMCNN processing.}
\label{fig:tsne_all}
\end{figure*}


\begin{figure*}[!htbp]
\centering
\includegraphics[width=0.32\linewidth]{C:/Users/devgu/Downloads/swati_hira_research_paper_work/final_kaggle_cells/results_ver_5/IndianPines_modulation_gamma.png} \hfill \includegraphics[width=0.32\linewidth]{C:/Users/devgu/Downloads/swati_hira_research_paper_work/final_kaggle_cells/results_ver_5/PaviaU_modulation_gamma.png} \hfill \includegraphics[width=0.32\linewidth]{C:/Users/devgu/Downloads/swati_hira_research_paper_work/final_kaggle_cells/results_ver_1/KSC_modulation_gamma.png}
\caption{Spatial Modulation Heatmaps ($\gamma$) for (a) Indian Pines, (b) PaviaU, and (c) KSC datasets. The heatmaps visualize the spatial scaling parameters learned by the FiLM layer, highlighting region-specific feature emphasis.}
\label{fig:gamma_all}
\end{figure*}

\begin{table*}[!htbp]
\centering
\caption{Noise Robustness Assessment (Indian Pines \& PaviaU)}
\label{tab:robustness}
\begin{tabular}{lcccc}
\toprule
\textbf{Dataset} & \textbf{Model} & \textbf{Clean OA (\%)} & \textbf{20dB SNR OA (\%)} & \textbf{10dB SNR OA (\%)} \\
\midrule
\multirow{2}{*}{Indian Pines} & Raw 2D-CNN & $86.42 \pm 0.65$ & $78.15 \pm 0.90$ & $68.40 \pm 1.25$ \\
& \textbf{SMCNN-MLP} & \textbf{$97.05 \pm 0.13$} & \textbf{$95.80 \pm 0.25$} & \textbf{$92.50 \pm 0.40$} \\
\midrule
\multirow{2}{*}{PaviaU} & Raw 2D-CNN & $92.15 \pm 0.40$ & $85.30 \pm 0.75$ & $76.10 \pm 1.10$ \\
& \textbf{SMCNN-MLP} & \textbf{$98.50 \pm 0.08$} & \textbf{$97.10 \pm 0.15$} & \textbf{$94.85 \pm 0.35$} \\
\bottomrule
\end{tabular}
\end{table*}





\begin{figure*}[!htbp]
\centering
\includegraphics[width=0.32\linewidth]{C:/Users/devgu/Downloads/swati_hira_research_paper_work/final_kaggle_cells/results_ver_6/confusion_matrix_IndianPines.png} \hfill \includegraphics[width=0.32\linewidth]{C:/Users/devgu/Downloads/swati_hira_research_paper_work/final_kaggle_cells/results_ver_6/confusion_matrix_PaviaU.png} \hfill \includegraphics[width=0.32\linewidth]{C:/Users/devgu/Downloads/swati_hira_research_paper_work/final_kaggle_cells/results_ver_6/confusion_matrix_KSC.png}
\caption{Confusion Matrices for (a) Indian Pines, (b) PaviaU, and (c) KSC datasets. Strong diagonal dominance visually aligns with the reported high accuracies.}
\label{fig:conf_all}
\end{figure*}


\begin{figure*}[!htbp]
\centering
\includegraphics[width=0.45\linewidth]{C:/Users/devgu/Downloads/swati_hira_research_paper_work/final_kaggle_cells/results_ver_last/classification_map_IndianPines.png} \hfill
\includegraphics[width=0.45\linewidth]{C:/Users/devgu/Downloads/swati_hira_research_paper_work/final_kaggle_cells/results_ver_last/classification_map_PaviaU.png} \\
\vspace{0.4cm}
\includegraphics[width=0.45\linewidth]{C:/Users/devgu/Downloads/swati_hira_research_paper_work/final_kaggle_cells/results_ver_last/classification_map_KSC.png}
\caption{Full Classification Maps for (a) Indian Pines, (b) PaviaU, and (c) KSC datasets. Sharp boundaries suggest effective AFDKF spatial edge preservation.}
\label{fig:class_all}
\end{figure*}




\subsection{Noise Robustness Validation}
To explicitly demonstrate the robustness provided by the AFDKF preprocessor, we evaluated the SMCNN-MLP against varying degrees of artificially injected Gaussian noise on both the Indian Pines and Pavia University datasets.

% Table X moved to appear with Figure 6 on float page 12

Table \ref{tab:robustness} demonstrates that while a standard 2D-CNN rapidly degrades under extreme noise conditions (68.40\% at 10dB), the SMCNN-MLP maintains a highly robust 92.50\% OA on Indian Pines, and 94.85\% on PaviaU, empirically supporting the robustness benefits of the proposed preprocessing pipeline.

\section{Limitations}
Audebert et al. (2019) similarly demonstrates that Salinas frequently reaches $\sim$99\% OA even under constrained scenarios, aligning with our observations. While the SMCNN-MLP achieves an excellent efficiency-accuracy trade-off, we explicitly acknowledge three limitations:
1) Performance degradation on highly mixed, low-resolution datasets (e.g., KSC) where computationally heavy models like HybridSN excel.
2) While MLP is shown to be equivalent to BiGRU under modulation, we have not yet directly compared against modern Mamba or Transformer-based FiLM modulators.
3) Edge deployment latency and actual watt-hour power draw on physical hardware (e.g., NVIDIA Jetson) remain untested.


While our framework establishes a strong efficiency benchmark against HybridSN and SSFTT, future work must explicitly evaluate lightweight MLP conditioning against emerging large-scale paradigms such as SpectralGPT. 

\subsection{Comparison with State Space Models}
Recent advancements in sequence modeling have introduced State Space Models (SSMs), notably MambaHSI and SpectralMamba, which achieve state-of-the-art accuracy in hyperspectral classification. While an experimental comparison would be illuminating, a direct empirical benchmarking against our proposed framework is not feasible within the scope of this study for three fundamental reasons. First, MambaHSI and SpectralMamba are large-scale architectures requiring extensive parameter counts and specialized GPU clusters, fundamentally conflicting with our explicit focus on "Green AI" and ultra-lightweight edge deployment. Second, SSMs rely on highly optimized, hardware-specific CUDA kernels (e.g., \texttt{causal-conv1d} and \texttt{mamba-ssm}) which lack universal support on low-power, resource-constrained edge devices such as UAV microcontrollers. Finally, the core theoretical advantage of Mamba lies in modeling ultra-long sequences with linear complexity; however, hyperspectral pixel signatures are relatively short (typically $\le 200$ bands). Our results demonstrate that for sequences of this length, simple affine modulation (FiLM) combined with robust preprocessing provides sufficient representational power at a fraction of the computational latency. Thus, while Mamba models represent the frontier of maximal accuracy, our framework operates on a strictly orthogonal Pareto frontier maximizing computational economy.

\section{Conclusion \& Future Work}\label{sec:conclusion}
We investigate whether robust preprocessing combined with lightweight FiLM modulation can challenge the necessity of increasingly complex HSI architectures, demonstrating that within the SMCNN framework, BiGRU does not consistently justify its computational overhead compared to lightweight modulation. By robustly denoising signals via an Adaptive Fast Desensitized Kalman Filter and extracting frequency patterns with a Polar-Logarithmic Continuous Wavelet Transform, we highlighted the critical role of data preparation over architectural complexity.

Our proposed framework, SMCNN, effectively utilizes Feature-wise Linear Modulation (FiLM) parameters to scale and shift standard 2D convolutions. Rigorous ablation studies revealed that our lightweight MLP variant matched the performance of highly complex BiGRU recurrent networks, establishing that lightweight modulation achieves a highly competitive accuracy-efficiency tradeoff. Ultimately, we present a lightweight FiLM-modulated CNN for HSI classification, preceded by Kalman-wavelet preprocessing, accompanied by a rigorous study showing that simple MLP modulation performs similarly to BiGRU modulation under strict training constraints, maximizing the accuracy-to-compute ratio. While the individual mathematical components (Kalman filters, wavelets, FiLM) are established techniques, their novel systems integration into an ultra-efficient pipeline demonstrates an alternative paradigm emphasizing preprocessing-driven efficiency. This establishes the SMCNN-MLP as a highly efficient, edge-deployment candidate for real-world UAV and satellite platforms. 

Future work will focus on deploying the SMCNN-MLP onto physical hardware, such as NVIDIA Jetson edge devices, to evaluate real-time inference latency and battery consumption under live agricultural and environmental monitoring conditions.



\begin{thebibliography}{99}
\bibitem{hypermamba2024} Liu, Q., Yue, J., Fang, Y., Xia, S., \& Fang, L. (2024). HyperMamba: A Spectral-Spatial Adaptive Mamba for Hyperspectral Image Classification. IEEE Transactions on Geoscience and Remote Sensing. \url{https://doi.org/10.1109/TGRS.2024.3482473}
\bibitem{scnet2025} Zhang, X., et al. (2025). SCNet: Spatial and Channel Reconstruction Convolutions for Lightweight Hyperspectral Image Classification. IEEE Geoscience and Remote Sensing Letters.
\bibitem{waveletdenoise2024} Li, Y., et al. (2024). Wavelet-guided deep learning for hyperspectral image denoising. ISPRS Journal of Photogrammetry and Remote Sensing.

\bibitem{yan2021} Li, R., Zheng, S., Duan, C., Yang, Y., \& Wang, X. (2020). Classification of hyperspectral image based on double-branch dual-attention mechanism network. Remote Sensing, 12(3), 582. \url{https://doi.org/10.3390/rs12030582}
\bibitem{kang2014} Kang, X., Li, S., \& Benediktsson, J. A. (2014). Spectral-spatial hyperspectral image classification with edge-preserving filtering. IEEE Transactions on Geoscience and Remote Sensing, 52(5), 2666-2677.
\bibitem{zhang2024} Huang, L., Chen, Y., \& He, X. (2024). Spectral-Spatial Mamba for Hyperspectral Image Classification. Remote Sensing, 16(13), 2449. \url{https://doi.org/10.3390/rs16132449}
\bibitem{lou2024} Lou, T.S., et al. (2024). Adaptive Fast Desensitized Kalman Filter for State Estimation. Circuits, Systems, and Signal Processing. \url{https://doi.org/10.1007/s00034-024-02801-3}
\bibitem{plcwt2023} Qian, Y., Ye, M., \& Zhou, J. (2013). Hyperspectral band selection exploiting linear prediction and wavelet filtering. IEEE Transactions on Geoscience and Remote Sensing, 51(7), 3680-3694.
\bibitem{melgani2004} Melgani, F., \& Bruzzone, L. (2004). Classification of hyperspectral remote sensing images with support vector machines. IEEE Transactions on Geoscience and Remote Sensing, 42(8), 1778-1790.
\bibitem{camps2005} Camps-Valls, G., \& Bruzzone, L. (2005). Kernel-based methods for hyperspectral image classification. IEEE Transactions on Geoscience and Remote Sensing, 43(6), 1351-1362.
\bibitem{fauvel2008} Fauvel, M., Benediktsson, J. A., Chanussot, J., \& Sveinsson, J. R. (2008). Spectral and spatial classification of hyperspectral data using SVMs and morphological profiles. IEEE Transactions on Geoscience and Remote Sensing, 46(11), 3804-3814.
\bibitem{bioucas2013} Bioucas-Dias, J. M., et al. (2013). Hyperspectral remote sensing data analysis and future challenges. IEEE Geoscience and Remote Sensing Magazine, 1(2), 6-36.
\bibitem{chen2014} Chen, Y., Lin, Z., Zhao, X., Wang, G., \& Gu, Y. (2014). Deep learning-based classification of hyperspectral data. IEEE Journal of Selected Topics in Applied Earth Observations and Remote Sensing, 7(6), 2094-2107.
\bibitem{makantasis2015} Makantasis, K., Karantzalos, K., Doulamis, A., \& Doulamis, N. (2015). Deep supervised learning for hyperspectral data classification through convolutional neural networks. IGARSS, 4959-4962.
\bibitem{zhao2016} Zhao, W., \& Du, S. (2016). Spectral-spatial feature extraction for hyperspectral image classification: A dimension reduction and deep learning approach. IEEE Transactions on Geoscience and Remote Sensing, 54(8), 4544-4554.
\bibitem{he2017} He, M., Li, B., \& Chen, H. (2017). Multi-scale 3D deep convolutional neural network for hyperspectral image classification. IEEE International Conference on Image Processing (ICIP), 3904-3908.
\bibitem{paoletti2018} Paoletti, M. E., Haut, J. M., Plaza, J., \& Plaza, A. (2018). A new deep convolutional neural network for fast hyperspectral image classification. ISPRS Journal of Photogrammetry and Remote Sensing, 145, 120-147.
\bibitem{li2019} Li, S., Song, W., Fang, L., Chen, Y., Ghamisi, P., \& Benediktsson, J. A. (2019). Deep learning for hyperspectral image classification: An overview. IEEE Transactions on Geoscience and Remote Sensing, 57(9), 6690-6709.
\bibitem{hong2020} Hong, D., Gao, L., Yao, J., Zhang, B., Plaza, A., \& Chanussot, J. (2021). Graph convolutional networks for hyperspectral image classification. IEEE Transactions on Geoscience and Remote Sensing, 59(7), 5966-5978.
\bibitem{roy2020} Roy, S. K., Krishna, G., Dubey, S. R., \& Chaudhuri, B. B. (2020). HybridSN: Exploring 3-D-2-D CNN feature hierarchy for hyperspectral image classification. IEEE Geoscience and Remote Sensing Letters, 17(2), 277-281.
\bibitem{sun2021} Sun, H., Zheng, X., Lu, X., \& Wu, S. (2020). Spectral-spatial attention network for hyperspectral image classification. IEEE Transactions on Geoscience and Remote Sensing, 58(5), 3232-3245.
\bibitem{zhong2018} Zhong, Z., Li, J., Luo, Z., \& Bell, M. (2018). Spectral-spatial residual network for hyperspectral image classification: A 3-D deep learning framework. IEEE Transactions on Geoscience and Remote Sensing, 56(1), 347-358.
\bibitem{hang2020} Hang, R., Li, Z., Ghamisi, P., Hong, D., Xia, G., \& Liu, Q. (2021). Classification of hyperspectral images via multitask generative adversarial networks. IEEE Transactions on Geoscience and Remote Sensing, 59(2), 1424-1436.
\bibitem{meng2021} Meng, Z., Li, L., Jiao, L., Feng, Z., Tang, X., \& Ding, M. (2019). Fully dense multiscale fusion network for hyperspectral image classification. IEEE Transactions on Geoscience and Remote Sensing, 57(11), 8911-8927.
\bibitem{wu2021} Wu, H., \& Prasad, S. (2018). Semi-supervised deep learning using pseudo-labels for hyperspectral image classification. IEEE Transactions on Image Processing, 27(3), 1259-1270.
\bibitem{ding2022} Ding, Y., Zhang, Z., Zhao, X., Cai, Y., Li, S., \& Yang, Y. (2021). Multi-scale receptive field graph attention network for hyperspectral image classification. IEEE Journal of Selected Topics in Applied Earth Observations and Remote Sensing, 15, 2307-2320.
\bibitem{liu2023} Liu, Y., \& Zheng, Y. (2023). Transformer-based spatial-spectral network for HSI classification. ISPRS Journal of Photogrammetry and Remote Sensing, 184, 1-15.
\bibitem{zhang2022} Zhang, X., \& Xu, J. (2020). 3D-CNN with channel attention for robust hyperspectral classification. Remote Sensing, 14(5), 1155.
\bibitem{chen2023} Ghamisi, P., et al. (2017). Advanced spectral-spatial classification of hyperspectral images. IEEE Geoscience and Remote Sensing Magazine, 5(1), 8-32.
\bibitem{gao2023} Gu, A., \& Dao, T. (2023). Mamba: Linear-time sequence modeling with selective state spaces. arXiv preprint arXiv:2312.00752.
\bibitem{huang2023} Huang, R., \& Zhao, Y. (2023). Fractional-order calculus in image denoising: Theory and applications. Signal Processing, 195, 108480.
\bibitem{schwartz2020} Schwartz, R., Dodge, J., Smith, N. A., \& Etzioni, O. (2020). Green AI. Communications of the ACM, 63(12), 54-63.

\bibitem{devguru2026repo} Tiwari, D., \& Hira, S. (2026). Source code for: Lightweight FiLM Modulation with Robust Preprocessing Challenges the Necessity of Complex Hyperspectral Architectures. GitHub Repository. Available online: \url{https://github.com/Devguru-codes/research_code}
\end{thebibliography}
\end{document}
